\chapter[\emj{🪣}~Counting, Radix y Bucket Sort (No Comparación)]{\emj{🪣}~Counting, Radix y Bucket Sort (Ordenar sin Comparar)}

\begin{dsaquees}

Todos los ordenamientos anteriores (Quick, Merge, Heap) COMPARAN pares de
elementos, y por eso ninguno baja de $O(n \log n)$. Estos tres hacen
trampa 🎯: no comparan, CUENTAN o REPARTEN, y logran $O(n)$ cuando los
datos cumplen ciertas condiciones (rango pequeño, enteros, distribución
uniforme).

\begin{itemize}
\item Counting Sort: cuenta cuántas veces aparece cada valor.
\item Radix Sort: ordena dígito por dígito.
\item Bucket Sort: reparte en cubetas y ordena cada una.
\end{itemize}

\end{dsaquees}

\begin{dsanino}

Imagina que quieres ordenar las cartas de un juego 🃏 numeradas del 1 al
10. En vez de comparar carta con carta, haces 10 montoncitos en la mesa,
uno por número. Avientas cada carta a su montoncito y luego recoges los
montoncitos en orden: 1, 2, 3\ldots ¡Listo, ya están ordenadas y nunca
comparaste dos cartas! Eso es Counting Sort.

Si los números fueran gigantes (como 375), lo haces por partes: primero
ordenas por el último dígito, luego por el de en medio, luego por el
primero. Como una fila que se reorganiza tres veces. Eso es Radix Sort.
Trucos que solo funcionan cuando los números no son "cualquier cosa".

\end{dsanino}

\begin{dsaotraforma}

Los ordenamientos por comparación tienen una cota inferior de
$\Omega(n \log n)$. Estos algoritmos la rompen porque \textbf{no comparan}
elementos entre sí; usan la estructura de los datos.

\textbf{Counting Sort}: cuenta la frecuencia de cada valor en un rango
$[0, k]$ y reconstruye el arreglo ordenado. $O(n + k)$. Estable. Solo
sirve si $k$ (rango) no es mucho mayor que $n$.

\textbf{Radix Sort}: aplica Counting Sort dígito por dígito, del menos al
más significativo (LSD). $O(d \cdot (n + b))$ con $d$ dígitos y base $b$.
Ideal para enteros o strings de longitud fija.

\textbf{Bucket Sort}: reparte los elementos en cubetas según su valor,
ordena cada cubeta (con otro sort) y concatena. $O(n)$ promedio si la
distribución es uniforme; $O(n^2)$ en el peor caso.

\end{dsaotraforma}

\begin{dsadiagrama}
\begin{verbatim}
COUNTING SORT de [2, 5, 3, 0, 2, 3, 0, 3]  (rango 0..5)

Conteo de cada valor:
  valor:  0  1  2  3  4  5
  veces:  2  0  2  3  0  1

Reconstruir en orden leyendo el conteo:
  0,0, 2,2, 3,3,3, 5  ->  [0, 0, 2, 2, 3, 3, 3, 5]

RADIX SORT de [170, 45, 75, 90, 802, 24, 2, 66]
  por unidades:  [170,90,802,2,24,45,75,66]
  por decenas :  [802,2,24,45,66,170,75,90]
  por centenas:  [2,24,45,66,75,90,170,802]  <- ordenado
\end{verbatim}
\end{dsadiagrama}

\begin{lstlisting}[language=C++]
COUNTING SORT (estable)
1. Cuenta cuántas veces aparece cada valor (arreglo count de tamaño k+1).
2. Acumula (prefix sum) para saber la posición final de cada valor.
3. Recorre el original de DERECHA a IZQUIERDA colocando cada elemento
   (así se mantiene estable).

📝 PSEUDOCÓDIGO (counting sort)
──────────────────────────────────────────────────────────────

función countingSort(a, k):        // valores en [0, k]
    count = arreglo de ceros de tamaño k+1
    PARA cada x en a: count[x]++
    PARA i desde 1 hasta k: count[i] += count[i-1]   // prefijos
    salida = arreglo del tamaño de a
    PARA i desde n-1 hasta 0:       // de derecha a izquierda = estable
        salida[--count[a[i]]] = a[i]
    return salida

💻 CÓDIGO C++ (Counting Sort)
──────────────────────────────────────────────────────────────

#include <vector>
using namespace std;

vector<int> countingSort(vector<int>& a, int k) {
    vector<int> count(k + 1, 0), out(a.size());
    for (int x : a) count[x]++;                 // frecuencias
    for (int i = 1; i <= k; i++) count[i] += count[i-1]; // prefijos
    for (int i = a.size() - 1; i >= 0; i--)     // estable
        out[--count[a[i]]] = a[i];
    return out;
}

💻 CÓDIGO C++ (Radix Sort, LSD para enteros >= 0)
──────────────────────────────────────────────────────────────

#include <algorithm>
void radixSort(vector<int>& a) {
    if (a.empty()) return;
    int maxv = *max_element(a.begin(), a.end());
    for (int exp = 1; maxv / exp > 0; exp *= 10) {  // dígito por dígito
        vector<int> out(a.size()), count(10, 0);
        for (int x : a) count[(x / exp) % 10]++;
        for (int i = 1; i < 10; i++) count[i] += count[i-1];
        for (int i = a.size() - 1; i >= 0; i--) {
            int d = (a[i] / exp) % 10;
            out[--count[d]] = a[i];
        }
        a = out;
    }
}

⏱️ COMPLEJIDAD (Big-O)
──────────────────────────────────────────────────────────────

  Algoritmo     │ Tiempo          │ Estable │ Cuándo usar
  ──────────────┼─────────────────┼─────────┼──────────────────
  Counting      │ O(n + k)        │ Sí      │ rango k pequeño
  Radix (LSD)   │ O(d·(n + b))    │ Sí      │ enteros/strings
  Bucket        │ O(n) promedio   │ Sí*     │ datos uniformes

* k = rango, d = número de dígitos, b = base.
\end{lstlisting}

\begin{dsaentrevistas}

✅ Rompen la cota O(n log n) porque NO comparan. Si un entrevistador
   pregunta "¿se puede ordenar en menos de O(n log n)?", la respuesta es
   "sí, si los datos lo permiten (rango acotado, enteros)".

💡 "Sort Colors" (LC 75): ordenar un arreglo de 0,1,2 es Counting Sort
   trivial (o Dutch flag en una pasada). "Maximum Gap" (LC 164) usa
   Bucket Sort para lograr O(n).

⚠️ Counting Sort explota si k es enorme: ordenar valores hasta $10^9$
   con counting necesitaría un arreglo de mil millones. Solo sirve con
   rango chico.

🔑 Estabilidad importa en Radix: cada pasada por dígito DEBE ser estable
   (por eso usa Counting Sort). Si no, el orden de los dígitos previos se
   pierde y el resultado es incorrecto.

\end{dsaentrevistas}

\begin{dsaresumen}

• No comparan: cuentan o reparten, rompen la cota O(n log n).
• Counting Sort: cuenta frecuencias; O(n + k); estable; rango k pequeño.
• Radix Sort: Counting dígito por dígito (LSD); enteros/strings.
• Bucket Sort: reparte en cubetas + ordena cada una; O(n) si uniforme.
• Radix depende de que cada pasada sea estable.
• Úsalos solo cuando los datos cumplan las condiciones (rango, tipo).
════════════════════════════════════════════════════════════════

\end{dsaresumen}
